\documentclass[11pt]{article}
\usepackage[T1]{fontenc}
\usepackage{textcomp}
\usepackage[margin=0.75in]{geometry}
\usepackage{amsmath,amssymb,amsthm}
\usepackage{array,booktabs}
\usepackage{hyperref}
\setlength{\parindent}{0pt}
\newtheorem{theorem}{Theorem}
\theoremstyle{definition}
\newtheorem{definition}{Definition}
\title{Flow Proof Artifact: prop-16-exterior-angle-greater}
\author{Generated by \texttt{flow doc proof}}
\date{Flow Proof Book}
\begin{document}
\maketitle
In a triangle, an exterior angle is greater than either remote interior angle.\\[0.5em]
\textbf{Source.} euclid — Elements, Book I, Proposition 16\\[0.5em]
\bigskip
\noindent{\large \textbf{Derived fact 1.} \textit{In a triangle, an exterior angle is greater than either remote interior angle} \hfill the Euclidean plane \cdot Euclid Book I \cdot Proposition 16: an exterior angle exceeds either remote interior angle}\par
\smallskip\noindent\textit{Coordinate: the Euclidean plane \cdot Euclid Book I \cdot Proposition 16: an exterior angle exceeds either remote interior angle}\par
\smallskip\noindent\textit{Source: euclid — Elements, Book I, Proposition 16}\par
\smallskip\noindent\textit{Built on: proposition 4: side-angle-side congruence, for Book I of the Elements in on the Euclidean plane, proposition 15: vertical angles are equal, for Book I of the Elements in on the Euclidean plane}\par
\medskip
\noindent\textbf{Goal.} In a triangle, an exterior angle is greater than either remote interior angle\par
\noindent$\displaystyle \angle ACD > \angle CAB$\par
\medskip
\noindent
\begin{tabular}{@{} >{\raggedright\arraybackslash}p{0.06\textwidth} >{\raggedright\arraybackslash}p{0.47\textwidth} >{\raggedleft\arraybackslash}p{0.41\textwidth} @{}}
\textbf{\#} & \textbf{Proof} & \textbf{Mathematics} \\
\midrule
\textbf{1.} & We prove that proposition 16: an exterior angle exceeds either remote interior angle for Book I of the Elements in the Euclidean plane. & \\[0.45em]
\textbf{2.} & Let D = point\_on\_BC\_produced\_beyond\_C. & \\[0.45em]
\textbf{3.} & Let F = midpoint\_of\_AC. & \\[0.45em]
\textbf{4.} & Let G = point\_on\_BF\_with\_FG\_equal\_to\_FA. & \\[0.45em]
\textbf{5.} & We invoke the derived fact: triangle\_ABC\_with\_exterior\_angle\_ACD. & \\[0.45em]
\textbf{6.} & We invoke the axiom governing Book I of the Elements in the Euclidean plane: proposition 4: side-angle-side congruence, for Book I of the Elements in on the Euclidean plane. & \\[0.45em]
\textbf{7.} & We invoke the derived fact governing Book I of the Elements in the Euclidean plane: proposition 15: vertical angles are equal, for Book I of the Elements in on the Euclidean plane. & \\[0.45em]
\textbf{8.} & From step 2, step 3, step 4, step 5, step 6, and step 7, by the chain of results established in Book I (step 2, step 3, step 4, step 5, step 6, and step 7), we obtain angle ACD > angle CAB. Hence proven. & $\displaystyle \angle ACD > \angle CAB$ \\[0.45em]
\bottomrule
\end{tabular}
\label{thm:Geometry-Euclid Book I-proposition_16_an_exterior_angle_exceeds_either_remote_interior_angle}
\par\medskip\hrule
\end{document}
